\section{Implementation}
\label{sec:implementation}

\subsection{Zephyr RTOS Integration}

The subsystem header (\texttt{drivers/biometrics.h}) defines the public API using static inline wrappers. Each public API function carries the \texttt{\_\_syscall} annotation, triggering kernel-mode dispatch code generation per entry point.

\begin{table}
\caption{Implementation Metrics (WeAct STM32F446RE)}
\label{tab:implementation_metrics}
\centering
\renewcommand{\arraystretch}{1.3}
\begin{tabularx}{\columnwidth}{|X|c|c|c|}
\hline
\textbf{Component} & \textbf{Flash (B)} & \textbf{Static RAM (B)} & \textbf{Heap (B)} \\
\hline
BASE Core API & 0 (Inlined) & 0 & 0 \\
\hline
ZFM-X0 Driver & 3,514 & 644 & 0 \\
\hline
GT-5X Driver & 4,288 & 144 & $\sim$520 \\
\hline
AI10 Driver & 4,974 & 2,176 & 0 \\
\hline
Typical Config & $<$ \SI{5}{\kilo\byte} & $<$ \SI{3}{\kilo\byte} & $<$ \SI{1}{\kilo\byte} \\
\hline
\end{tabularx}
\vspace{2pt}
\footnotesize{\\ \textit{Note: GT-5X uses dynamic allocation for packet buffers; ZFM-X0 and AI10 use static buffers. AI10 RAM includes dual frame buffers and continuous match thread state. AI10 figures reflect a pre-upstream implementation; the driver is under active development and footprint may change prior to merge.}}
\end{table}

\subsection{Driver Implementations}

\subsubsection{ZFM-X0 Family Driver}

The ZFM-X0 driver targets the ZhianTec protocol family (ZFM-20/60/70 and rebrands including Adafruit, Grow R307, and Hi-Link capacitive modules). It uses a variable-length packet format with checksum validation and fully interrupt-driven UART communication via semaphores. The driver enforces the four-state enrollment machine and supports both optical and capacitive sensors with identical code.

\subsubsection{GT-5X Family Driver}

The GT-5X driver targets the ADH-Tech optical family (GT-511C3/521F32/521F52). It implements the fixed 12-byte command protocol with 29 opcodes. The driver uses dynamic allocation only for packet buffers while satisfying the $O(1)$ memory contract and requires three samples for enrollment.

\subsubsection{DFRobot AI10 Driver (Multi-Modal)}

The AI10 driver targets the DFRobot SEN0677 module, a binocular camera-based sensor supporting face recognition and palm vein detection on a single device. This is the first driver in the BASE ecosystem to demonstrate the multi-modal capability: \texttt{get\_capabilities()} reports \texttt{supported\_modalities = BIOMETRIC\_MODALITY\_FACE | BIOMETRIC\_MODALITY\_PALM}. The driver maps face UIDs (1--1000) and palm UIDs (1001--2000) transparently into the standard \texttt{biometric\_match\_result} structure by setting the \texttt{modality} field accordingly. Enrollment uses a single-shot command, but the application loop remains identical across all drivers because it checks \texttt{samples\_required} (set to 1 for AI10).

\subsubsection{Protocol Comparison}

Table~\ref{tab:protocol_comparison} summarizes the three driver families.

\begin{table*}[!ht]
\caption{Driver Protocol Comparison}
\label{tab:protocol_comparison}
\centering
\renewcommand{\arraystretch}{1.2}
\setlength{\tabcolsep}{4pt}
\begin{tabularx}{\textwidth}{|l|X|X|X|}
\hline
\textbf{Feature} & \textbf{ZFM-X0} & \textbf{GT-5X} & \textbf{DFRobot AI10} \\
\hline
Sensor Type & Optical/Capacitive contact & Optical contact & Camera-based, binocular \\
\hline
Modalities & Fingerprint & Fingerprint & Face + Palm vein \\
\hline
Start Code & \texttt{0xEF01} & \texttt{0x55AA}/\texttt{0x5AA5} & \texttt{0xEF 0xAA} \\
\hline
Packet Size & Variable (9--267 B) & Fixed (12 B) & Variable + XOR \\
\hline
Enrollment Samples & 2 & 3 & 1 (single-shot) \\
\hline
Checksum & Payload sum & All-bytes sum & XOR over header+data \\
\hline
Memory & Static buffers & Dynamic heap & Static buffers \\
\hline
\texttt{supported\_modalities} & \texttt{FINGERPRINT} & \texttt{FINGERPRINT} & \texttt{FACE | PALM} \\
\hline
\end{tabularx}
\smallskip
\par\noindent\footnotesize\textit{Note: ZFM-X0 supports both optical and capacitive sensors with identical driver code. AI10 UID range 1--1000 maps to face; 1001--2000 maps to palm, handled transparently by the driver.}
\end{table*}

\subsection{Devicetree Integration}

BASE drivers use Zephyr's devicetree for hardware configuration. Example binding for the ZFM-X0 family:

\begin{small}
\begin{verbatim}
&uart2 {
  status = "okay";
  current-speed = <57600>;
  fingerprint0: zfm60@0 {
    compatible = "zhiantec,zfm-x0";
    comm-addr = <0xFFFFFFFF>;
    password = <0x00000000>;
  };
};
\end{verbatim}
\end{small}

We use devicetree properties to configure communication parameters like baud rate, address, and other sensor settings. \texttt{DT\_INST\_FOREACH\_STATUS\_OKAY} ensures automatic device structure generation. This allows multiple sensors to coexist without manual registration code.

\subsection{Robustness and Validation}
All blocking operations accept \texttt{k\_timeout\_t} parameters. \texttt{K\_NO\_WAIT} can be used for instant operations, \texttt{K\_FOREVER} can be used for indefinite wait, or specific timeouts for deterministic timeouts. Overflow protection is capped at 3600 seconds. Template IDs are bounds-checked against \texttt{max\_templates}, attribute values are range-checked, and timeout values are clamped. UART RX handlers detect overflow, invalid lengths, and checksum mismatches. It sets error flags that are checked after semaphore acquisition to prevent invalid data.

We validated the drivers on six physical devices (GT-511C3, DY-50, JM-101, Grow R307, HLK-ZW3021, HLK-ZW0623) and the WeAct STM32F446RE board, confirming cross-architecture compatibility and protocol-level abstraction across optical and capacitive sensors.

\subsection{Upstream Availability}
\label{sec:availability}
The reference implementation of BASE has been merged\cite{zephyr_pr, zephyr_pr2} into the Zephyr upstream kernel and released as part of Zephyr 4.4\cite{zephyr_44_notes}. As the designated maintainer of this subsystem, the first author ensures ongoing architectural consistency and the integration of future biometric modalities for the open-source industrial community.